
\section*{Introduction}

SOD2 proteins protect cells against oxidative damage by converting toxic
superoxide radicals into hydrogen peroxide and diatomic oxygen. SOD2 proteins
can bind either manganese (MnSOD) or iron (FeSOD) as co-factors; their
structure consists of conserved residues His26, His74, Asp159, and His163
(human sequence coordinates) forming the anchors of a trigonal bypyramidal
metal pocket and variations in the outer shell, based on bound metal, to ensure
superoxide dismutase activity \cite{perry2010bbapap}. SOD2 is present
throughout tree of life, with bacterial SOD2 often preferring Fe as cofactor,
mammalian SOD2 almost always requires Mn as a cofactor to function in
mitochondria, and metal preference not always phylogenetically constrained
\cite{sendra2023nee}.

Existing literature is vague on the types of SOD2 proteins in coral hosts and
their symbiotic algae. While there is consensus on the presence of
mitochondrial MnSOD in corals, other forms of SOD2 in coral and symbiotic algae
are rarely discussed/analyzed separately. \cite{saragosti2010plosone} reports
experimental evidence for secreted superoxide dismutase activity in
\textit{Stylophora pistillata} without identifying the source of such activity.
\cite{dash2007gene} experimentally identifies two forms of MnSOD in
\textit{Hydra vulgaris}, a non-coral species in the cnidaria phylum.
\cite{krueger2015eb} computationally compares various annotated MnSOD and FeSOD
sequences from various symbiodinium species, with grouping of sequences, based
on putative annotation, inconsistent with findings of this paper. Indeed, most
transcriptomic analyses in the field depend on blasting against UniProt or
NCBI, or matching against the HMM profile K04564, a pan-metagenomic KEGG
ortholog (\cite{kegg}) for SOD2. Neither approacher provide resolution beyond
putative superoxide dismutase activity.

This paper presents results from an unbiased computational search and grouping
of SOD2 protein sequences, using NCBI deposited genome sequences from 36 coral
and 10 symbiodinium species, as well as RNAseq datasets from 5 papers.

This work here looked at X genomes and Y transcriptomes to classify SOD2
proteins. When protein sequences are available, they are used. In all cases,
6-frame translations of the genome sequences were searched exhaustively against
KO database to identify likely proteins. Together the curated protein sequences
and HMM searched protein sequences form the search database. We then started
with K04564 match regions in these proteins, filtered away those w/o key
residues that are necessary for SOD2 to form a metal binding pocket. For each
remaining sequence, we searched for possible leader sequences up to 75 AAs
upstream of the start of the PF00081, using DeepLoc. The sequences were then
grouped by leader pattern, within coral hosts and within algae, separately.
